\uuid{pCpP}
\titre{Nombres premiers}
\auteur{bodin ; croizat ; sacre}
\organisation{exo7}
\licence{CC BY-NC-SA 4.0}
\datecreate{2022-03-08}
\niveau{L1}
\module{Algèbre}
\chapitre{Arithmétique dans Z}
\sousChapitre{Nombres premiers, nombres premiers entre eux}
\difficulte{3}
\theme{Nombres premiers}
\exo7qcmid{hilisit.arithmetique.036}

\contenu{
\begin{question}
\qclassification{103.04}
    Les entiers suivants ont été factorisés correctement.
    Quelles sont les écritures qui sont des décompositions en facteurs premiers ?
    \begin{answers}
        \bad{$3\,025 = 1^3 \times 5^2 \times 11^2$}
        \bad{$1\,836 = 2^2 \times 3 \times 3^2 \times 17$}
        \bad{$1\,444\,716 = 2^2 \times 7^3 \times 9^2 \times13$}
        \good{$13\,915 = 5 \times 11^2 \times 23$}
    \end{answers}
    \begin{explanations}
    Chaque facteur doit être de la forme $p_i^{\alpha_i}$ avec $p_i$ un nombre premier (donc pas $1$ ni $9$) et $\alpha_i>0$. En plus les $p_i$ doivent être deux à deux distincts. Avec ces contraintes la décomposition est unique (à l'ordre des facteurs près).
    \end{explanations}
\end{question}
}
